% ---------------------------------------------------------------------------
%  Relatório da Semana 3 — Tabuleiro de Xadrez Eletrônico
%
%  Compilar:  make pdf3        (usa xelatex, dentro do devShell do flake)
%  ou:        xelatex -output-directory=docs docs/Relatorio-3.tex   (2 passadas)
%
%  As figuras vêm de docs/diagramas/*.png, geradas dos .mmd por `mmdc`
%  (alvo `make diagramas`).
% ---------------------------------------------------------------------------
\documentclass[12pt,a4paper]{article}

\usepackage{fontspec}
\usepackage{polyglossia}
\setdefaultlanguage{portuguese}

\setmainfont{DejaVu Sans}
\setmonofont{DejaVu Sans Mono}[Scale=0.85]

\usepackage[margin=2.5cm]{geometry}
\usepackage{graphicx}
% Vale tanto para `xelatex docs/Relatorio-3.tex` da raiz (o que `make pdf3`
% faz, por causa do -output-directory) quanto de dentro de docs/.
\graphicspath{{docs/diagramas/}{diagramas/}}
\usepackage{float}
\usepackage{amssymb}       % \checkmark nas tabelas de rastreabilidade
\usepackage{booktabs}
\usepackage{longtable}
\usepackage{array}
\usepackage{xcolor}
\usepackage{fancyvrb}
\usepackage{enumitem}
\usepackage[hidelinks]{hyperref}

\floatplacement{figure}{H}
\setlist{itemsep=2pt, parsep=0pt, topsep=4pt}

% Colunas de largura fixa sem justificação: as células das tabelas são curtas
% e a justificação nelas só produz hifenização feia e espaços largos.
\newcolumntype{L}[1]{>{\raggedright\arraybackslash}p{#1}}

% Blocos de saída de terminal / protocolo.
\definecolor{codebg}{HTML}{F4F4F6}
\DefineVerbatimEnvironment{terminal}{Verbatim}{%
  fontsize=\footnotesize, frame=leftline, framerule=1pt,
  rulecolor=\color{gray}, framesep=8pt, xleftmargin=6pt}

\newcommand{\diagrama}[3]{%
  \begin{figure}[H]
    \centering
    \includegraphics[width=#2\textwidth]{#1}
    \caption{#3}
    \label{fig:#1}
  \end{figure}}

\title{Tabuleiro de Xadrez Eletrônico com Integração a \emph{Chess Engine}\\[6pt]
       \large Relatório da Semana 3}
\author{Grupo W\\ PCS3732 --- Laboratório de Processadores}
\date{Agosto de 2026}

\begin{document}

\maketitle

\begin{abstract}
\noindent
A terceira semana do projeto foi dedicada à camada de baixo nível e à interface
com o usuário. O processo em C foi reescrito de forma modular e ganhou uma
segunda camada de entrada --- um teclado matricial 4$\times$4 em que os lances
são digitados ---, intercambiável com a matriz de \emph{reed switches} e
falando exatamente o mesmo protocolo IPC com a aplicação Python. Esse é o
\textbf{plano B} do projeto: a montagem do tabuleiro instrumentado avançou, mas
não deve ficar pronta no prazo, e com o teclado o sistema inteiro (Stockfish,
Lichess, GUI, roque, capturas) continua jogável. Em paralelo, as escolhas que
antes só existiam na linha de comando passaram a ter interface gráfica, e a
suíte automatizada ganhou 30 verificações novas que exercitam o binário em C de
ponta a ponta.
\end{abstract}

\tableofcontents
\newpage

% ===========================================================================
\section{Desenvolvimento de Software}
% ===========================================================================

\subsection{Reestruturação do processo em C}

Na semana 2 o código de baixo nível era uma adaptação direta da prova de
conceito feita para Arduino: um único arquivo que inicializava a
\texttt{wiringPi}, varria a matriz e imprimia as mudanças. Nesta semana ele foi
reescrito como um programa modular, com responsabilidades separadas em
arquivos próprios e uma interface de linha de comando (Tabela~\ref{tab:modulos}).

\begin{table}[H]
\centering
\caption{Módulos do processo em C (\texttt{src/}).}
\label{tab:modulos}
\begin{tabular}{@{}ll@{}}
\toprule
\textbf{Arquivo} & \textbf{Responsabilidade} \\
\midrule
\texttt{main.c}         & CLI: escolhe a camada de entrada e liga-a ao canal IPC \\
\texttt{reed\_layer.c}  & Camada 1: varredura da matriz 8$\times$8 e \emph{debouncing} \\
\texttt{keypad\_layer.c}& Camada 2: teclado matricial 4$\times$4, lances digitados \\
\texttt{board.c}        & Casas do tabuleiro e espelho de ocupação \\
\texttt{ipc.c}          & Serialização dos eventos (\emph{stdout} ou \emph{Named Pipe}) \\
\texttt{lcd.c}          & Display 16$\times$2 por I\textsuperscript{2}C (opcional) \\
\texttt{gpio.c}         & \emph{Wrapper} da \texttt{wiringPi} (compila também sem ela) \\
\texttt{runstate.c}     & Encerramento limpo em \texttt{SIGINT}/\texttt{SIGTERM} \\
\bottomrule
\end{tabular}
\end{table}

A decisão de projeto que organiza tudo isso é a de \textbf{camadas de entrada
intercambiáveis}: as duas camadas emitem exatamente os mesmos eventos, então a
aplicação Python não sabe --- nem precisa saber --- qual delas está do outro
lado. Trocar de uma para a outra é trocar uma opção de linha de comando
(\texttt{-{}-input reed} ou \texttt{-{}-input keypad}), sem tocar em nada da
lógica de jogo.

\diagrama{arq_c.png}{0.70}{Arquitetura do processo em C. As duas camadas de
entrada convergem para o mesmo módulo de IPC, o que as torna intercambiáveis do
ponto de vista da aplicação Python.}

Outro ponto de projeto é a detecção da \texttt{wiringPi} em tempo de
compilação. A biblioteca só existe no Raspberry Pi; sem ela o binário compila
do mesmo jeito (as funções de GPIO viram \emph{stubs} que recusam a execução),
o que permite exercitar o \emph{parser} do teclado em qualquer máquina --- e é
exatamente disso que a suíte de testes da Seção~\ref{sec:testes} se aproveita.

\subsection{Plano B: camada de teclado matricial}

O risco identificado na semana 2 --- a chapa de acrílico comprada era grande
demais para o alcance dos ímãs, e a remontagem do tabuleiro consome tempo ---
motivou uma alternativa de contingência: em vez de mover a peça no tabuleiro
instrumentado, o jogador \textbf{digita o lance} num teclado matricial 4$\times$4
de baixo custo. Todo o resto do sistema permanece igual.

\paragraph{Impacto nos requisitos.} O plano B não é neutro do ponto de vista da
especificação, e vale declará-lo: com o teclado, \textbf{RF1} (``detecção das
peças no tabuleiro físico e de seus movimentos'') deixa de ser satisfeito ---
o lance é informado, não detectado --- e \textbf{RNF3} (robustez contra falsos
positivos por esbarrões) perde o objeto, já que não há sensor a ser
perturbado. \textbf{RF2}, \textbf{RF3} e \textbf{RNF2} continuam valendo
integralmente, e são exatamente eles que o teclado preserva. A camada de
\emph{reed switches} continua implementada e é o caminho principal: o teclado é
contingência para o caso de a montagem não fechar, não uma troca de escopo. A
Seção~\ref{sec:rastreabilidade} registra esse estado requisito a requisito.

O obstáculo é que o teclado tem apenas as letras \texttt{A}--\texttt{D} e o
tabuleiro precisa de \texttt{a}--\texttt{h}. A solução adotada é o
\textbf{toque repetido}: a mesma tecla pressionada duas vezes seguidas vale
pela letra do bloco seguinte (\texttt{A}$\to$\texttt{a}, \texttt{AA}$\to$\texttt{e}),
e um terceiro toque volta ao início, o que permite corrigir sem apagar. Não há
temporizador envolvido --- como o lance alterna coluna e fileira, duas letras
seguidas só podem ser a mesma casa sendo redigitada, o que torna a
desambiguação determinística.

\begin{table}[H]
\centering
\caption{Exemplos de digitação de lances.}
\label{tab:teclas}
\begin{tabular}{@{}lll@{}}
\toprule
\textbf{Lance} & \textbf{Teclas} & \textbf{Evento emitido} \\
\midrule
\texttt{a2a4} & \texttt{A 2 A 4 \#}       & \texttt{a2:0,a4:1} \\
\texttt{e2e4} & \texttt{AA 2 AA 4 \#}     & \texttt{e2:0,e4:1} \\
\texttt{g1f3} & \texttt{CC 1 BB 3 \#}     & \texttt{g1:0,f3:1} \\
Roque curto   & \texttt{AA 1 CC 1 \#} e depois \texttt{DD 1 BB 1 \#} & \texttt{e1:0,g1:1} e \texttt{h1:0,f1:1} \\
\bottomrule
\end{tabular}
\end{table}

A tecla \texttt{\#} confirma o lance e a \texttt{*} apaga a última tecla. O
roque é digitado em dois lances (rei e depois torre), exatamente como já
acontece no tabuleiro físico, o que reaproveita sem alteração o tratamento de
roque em duas etapas já existente no Python.

\diagrama{fluxo_keypad.png}{0.92}{Fluxo de digitação de um lance na camada de
teclado, incluindo os comandos e a recusa local.}

O processo em C mantém um \textbf{espelho} de onde estão as peças do jogador e
recusa na origem o que o tabuleiro físico também recusaria: mover a partir de
uma casa vazia, mover para uma casa já ocupada por peça própria, ou origem
igual ao destino. A legalidade do lance continua sendo decidida pelo Python
--- essa validação local só evita gerar eventos que o tabuleiro real jamais
produziria. Como o espelho pode divergir do tabuleiro virtual (por exemplo,
quando o oponente captura uma peça), há um conjunto de comandos com prefixo
\texttt{0}:

\begin{table}[H]
\centering
\caption{Comandos do teclado (equivalentes a mexer numa peça sem fazer um lance).}
\begin{tabular}{@{}lll@{}}
\toprule
\textbf{Comando} & \textbf{Efeito} & \textbf{Quando usar} \\
\midrule
\texttt{0 1} \emph{casa} \texttt{\#} & Retira a peça da casa & O oponente capturou uma peça \\
\texttt{0 2} \emph{casa} \texttt{\#} & Coloca uma peça na casa & A aplicação pediu para repor uma peça \\
\texttt{0 9 \#} & Reenvia o estado das 64 casas & Espelho e tabuleiro virtual divergiram \\
\texttt{0 0 \#} & Volta o espelho à posição inicial & Recomeçar do zero \\
\bottomrule
\end{tabular}
\end{table}

\subsection{Extensão do protocolo IPC: linhas \texttt{@entry}}

Na camada de \emph{reed switches}, a peça na mão do jogador é visível pelos
próprios sensores, e a GUI usa essa informação para destacar os destinos
legais. Digitando o lance não há sensor nenhum --- e sem alguma informação
adicional o jogador digitaria às cegas até o \texttt{\#}.

A solução foi estender o protocolo IPC com um segundo tipo de linha, que
começa com \texttt{@} --- caractere que nenhuma casa do tabuleiro pode ter, de
modo que qualquer leitor que só entenda \texttt{casa:estado} as ignora sem
tropeçar:

\begin{terminal}
@entry|origem|destino|texto|status
\end{terminal}

\noindent
Digitando \texttt{AA2AA4\#} (lance \texttt{e2e4}), a sequência emitida é:

\begin{terminal}
@entry|||Lance: A_|Digitando...
@entry|||Lance: E_|Coluna trocada
@entry|e2||Lance: E2_|Digitando...
@entry|e2||Lance: E2A_|Digitando...
@entry|e2||Lance: E2E_|Coluna trocada
@entry|e2|e4|Lance: E2E4|Digitando...
e2:0,e4:1
@entry||||Enviado e2e4
\end{terminal}

Com isso a GUI passa a destacar a casa de origem e os destinos legais da peça
escolhida assim que as duas primeiras teclas formam uma casa, e a barra de
status mostra o que já foi digitado. O destino digitado, ainda não confirmado,
ganha uma borda âmbar. Só peças do próprio jogador são destacadas: apontar uma
casa vazia ou uma peça do oponente é erro de digitação, e destacá-la sugeriria
um lance inexistente.

\subsection{Interface de configuração}
\label{sec:menu}

Na revisão por pares da semana 2, o Grupo I sugeriu substituir a configuração
por CLI por uma interface visual interativa, e o Grupo B pediu melhor
sinalização de erros ao usuário. Ambas as sugestões foram atendidas nesta
semana.

Tudo o que os alvos do \texttt{Makefile} escolhem por variáveis --- modo de
jogo, oponente, cor das peças, camada de entrada, tempo da \emph{engine},
controle de tempo do Lichess --- passou a ser escolhido também na própria
interface, descrito num único módulo (\texttt{app/launcher.py}) e renderizado
de duas formas: uma janela \emph{pygame} operada por teclado ou mouse, e uma
lista numerada no terminal para quando não há \emph{display} (acesso por SSH).
Executar \texttt{python -m app.main} sem argumento nenhum abre o menu.

\begin{terminal}
Tabuleiro de Xadrez Eletrônico
Stockfish (1s por lance) - peças brancas - entrada: mock (gui)
--------------------------------------------------------------
  Modo de jogo                          Stockfish local
  Tempo por lance (s)              <       1.0       >
  Cor das peças físicas                          brancas
  Entrada do tabuleiro               mock (sem hardware)
  ...
  INICIAR PARTIDA
  Compilar processo C (make board-input)
  Sair
--------------------------------------------------------------
Equivalente no terminal: make stockfish COLOR=white
Setas: escolher/alterar - Enter: editar ou iniciar - F5: compilar
\end{terminal}

Quanto à sinalização de erros, o menu \textbf{verifica as pré-condições antes
de começar} e impede a partida enquanto houver impedimento: token do Lichess
ausente, controle de tempo que a \emph{Board API} recusa, processo em C ainda
não compilado, Stockfish fora do \texttt{PATH}. Cada um vira uma mensagem
explicativa no rodapé, em vez de um erro no meio da partida. O menu também
compila o processo em C sem sair da interface e exibe o comando equivalente
(\texttt{make ...}), de modo que a CLI continua disponível para quem
desenvolve.

\subsection{Ações durante a partida}

Com a partida em andamento, a tecla \texttt{Esc} passou a abrir um menu de
ações sobre o tabuleiro: reiniciar (modo Stockfish), abortar, oferecer ou
aceitar empate, desistir, voltar ao menu de configuração ou sair. Antes, uma
proposta de empate do oponente só podia ser respondida pelo site do Lichess.

Duas decisões de interface merecem registro. Primeiro, enquanto a partida está
viva toda ação irreversível pede confirmação --- num tabuleiro operado por
toque, um encostão não pode custar a partida. Segundo, ao fim da partida o
menu abre sozinho e sem confirmação, já com as opções que ainda fazem sentido,
e ``Voltar ao menu'' devolve o jogador à configuração sem reabrir o programa.

% ===========================================================================
\section{Desenvolvimento de Hardware}
% ===========================================================================

A montagem do tabuleiro instrumentado avançou: cerca de \textbf{duas fileiras
completas} de casas foram montadas com \emph{reed switches} e diodos, o
suficiente para confirmar em escala real o mapeamento de linhas e colunas e a
detecção através da chapa. Ainda assim, a avaliação do grupo é que
\textbf{não haverá tempo de concluir as 64 casas} dentro do cronograma da
disciplina --- o que confirma o risco levantado na semana 2 e justifica o plano
B descrito acima.

Do lado do teclado matricial, a implementação foi \textbf{testada com o
hardware de verdade}, reaproveitando a fiação já montada na bancada (a mesma do
experimento 6): linhas nos pinos BCM \texttt{16, 20, 21, 26} e colunas em
\texttt{19, 13, 6, 5}. Para essa verificação o processo em C ganhou um modo de
conferência de bancada:

\begin{terminal}
$ ./build/board_input --input keypad --raw
[bancada] tecla 'A'  (linha 0 = pino 16, coluna 3 = pino 19)
\end{terminal}

\noindent
Nesse modo nenhum evento é enviado --- ele só mostra em que interseção da
matriz cada tecla foi lida, o que separa um problema de fiação de um problema
de lógica. A conferência confirmou que o teclado da bancada é o 4$\times$4
padrão --- \texttt{A B C D} na quarta coluna, \texttt{* 0 \# D} na última
linha ---, que é exatamente o mapeamento assumido em
\texttt{keypad\_config\_default()}: nenhuma tradução de rótulos foi
necessária.

A matriz de \emph{reed switches} usa pinos distintos (linhas
\texttt{4, 5, 6, 12, 13, 16, 19, 20}, colunas \texttt{21, 22, 23, 24, 25, 26,
27, 17}), de modo que as duas camadas podem coexistir fisicamente no mesmo
Raspberry Pi.

% ===========================================================================
\section{Evidências de Testes Experimentais}
\label{sec:testes}
% ===========================================================================

\subsection{Suíte automatizada}

A suíte de testes ganhou uma quarta \emph{suite},
\texttt{tests/test\_keypad\_layer.py}, com \textbf{30 verificações}. Ela não é
um teste de unidade do Python: ela executa o binário em C de verdade
(\texttt{build/board\_input --input keypad --keys stdin}), digita as teclas na
entrada padrão e confere os eventos que saem --- ou seja, exercita o mesmo
caminho de código do teclado real, trocando apenas a origem das teclas. Além
disso, as linhas \texttt{@entry} são conferidas com o \emph{parser} que a
própria aplicação usa (\texttt{app.ipc\_reader.parse\_entry}), o que valida os
dois lados do protocolo. A suíte é pulada, sem falhar, quando o binário não
foi compilado.

Resultado da execução completa (\texttt{make test}):

\begin{terminal}
======================================================================
Todas as 4 suítes passaram.
\end{terminal}

\subsection{Verificações da camada de teclado}

A Tabela~\ref{tab:evidencias} reproduz as verificações executadas, agrupadas
por objetivo. Todas passaram.

\begin{longtable}{@{}L{0.32\textwidth}L{0.28\textwidth}L{0.32\textwidth}@{}}
\caption{Evidências da camada de teclado (\texttt{tests/test\_keypad\_layer.py}).}
\label{tab:evidencias}\\
\toprule
\textbf{Verificação} & \textbf{Teclas} & \textbf{Saída observada} \\
\midrule
\endfirsthead
\toprule
\textbf{Verificação} & \textbf{Teclas} & \textbf{Saída observada} \\
\midrule
\endhead
\multicolumn{3}{@{}l}{\textit{Colunas: tecla simples e tecla repetida}} \\[2pt]
\texttt{a2a4} com teclas simples & \texttt{A2A4\#} & \texttt{a2:0,a4:1} \\
\texttt{e2e4} com tecla repetida & \texttt{AA2AA4\#} & \texttt{e2:0,e4:1} \\
\texttt{b1c3} (cavalo) & \texttt{B1C3\#} & \texttt{b1:0,c3:1} \\
\texttt{g1f3} (cavalo) & \texttt{CC1BB3\#} & \texttt{g1:0,f3:1} \\
\texttt{d2h6} mistura os blocos & \texttt{D2DD6\#} & \texttt{d2:0,h6:1} \\
Terceiro toque volta a \texttt{a} & \texttt{AAA2AAA4\#} & \texttt{a2:0,a4:1} \\
\addlinespace
\multicolumn{3}{@{}l}{\textit{Correção com \texttt{*}}} \\[2pt]
Apaga a fileira errada & \texttt{AA2AA3*4\#} & \texttt{e2:0,e4:1} \\
Apaga tudo e redigita & \texttt{AA2****A2A4\#} & \texttt{a2:0,a4:1} \\
\texttt{\#} com entrada incompleta & \texttt{AA2\#} & nada emitido \\
\addlinespace
\multicolumn{3}{@{}l}{\textit{Recusas locais (o que o tabuleiro físico também recusaria)}} \\[2pt]
Origem vazia & \texttt{AA5AA6\#} & nada emitido \\
Destino ocupado & \texttt{A1A2\#} & nada emitido \\
Origem igual ao destino & \texttt{AA2AA2\#} & nada emitido \\
\addlinespace
\multicolumn{3}{@{}l}{\textit{Comandos}} \\[2pt]
\texttt{0-1} retira a peça & \texttt{01AA2\#} & \texttt{e2:0} \\
\texttt{0-2} coloca a peça & \texttt{02AA5\#} & \texttt{e5:1} \\
\texttt{0-1} em casa já vazia & \texttt{01AA5\#} & nada emitido \\
\texttt{0-2} em casa já ocupada & \texttt{02AA2\#} & nada emitido \\
\addlinespace
\multicolumn{3}{@{}l}{\textit{Sequências de partida}} \\[2pt]
Espelho acompanha os lances & \texttt{AA2AA4\#AA4AA5\#} & \texttt{e2:0,e4:1} e \texttt{e4:0,e5:1} \\
Roque curto em dois lances & \texttt{AA2AA4\# CC1BB3\# BB1C4\# AA1CC1\# DD1BB1\#} & \texttt{e1:0,g1:1} seguido de \texttt{h1:0,f1:1} \\
De pretas, \texttt{e7e5} & \texttt{AA7AA5\#} & \texttt{e7:0,e5:1} \\
De pretas, \texttt{e2} está vazia & \texttt{AA2AA4\#} & nada emitido \\
\addlinespace
\multicolumn{3}{@{}l}{\textit{Reenvio do estado completo}} \\[2pt]
\texttt{0-9} reenvia as 64 casas & \texttt{09\#} & 1 linha com as 64 casas \\
Fileiras 1 e 2 ocupadas & \texttt{09\#} & estado inicial correto \\
\addlinespace
\multicolumn{3}{@{}l}{\textit{Digitação em andamento (linhas \texttt{@entry})}} \\[2pt]
Cada tecla emite uma linha & \texttt{AA2A} & 4 linhas \texttt{@entry} \\
Origem sai com a casa completa & \texttt{AA2} & \texttt{origin='e2'}, \texttt{text='Lance: E2\_'} \\
Casa incompleta não vira origem & \texttt{AA2A} & \texttt{origin='e2'} (inalterada) \\
Destino sai na quarta tecla & \texttt{AA2AA4} & \texttt{target='e4'} \\
\texttt{\#} limpa a digitação & \texttt{AA2AA4\#} & \texttt{status='Enviado e2e4'} \\
Comando não emite origem & \texttt{01AA2} & \texttt{origin=None} \\
\texttt{*} desfaz a origem & \texttt{AA2*} & \texttt{status='Apagado'} \\
Linhas \texttt{@entry} não contaminam os eventos & --- & canais separados \\
\bottomrule
\end{longtable}

\subsection{O que essas evidências cobrem}

Vale destacar três resultados que respondem a riscos concretos do projeto:

\begin{itemize}
  \item \textbf{O roque funciona pelo teclado.} A sequência de cinco lances
        produz exatamente os mesmos eventos que o tabuleiro físico produziria
        em duas etapas --- ou seja, o plano B não perde nenhuma regra especial
        do xadrez.
  \item \textbf{A troca de camada é transparente.} Digitar \texttt{AA2AA4\#}
        gera \texttt{e2:0,e4:1}, exatamente o que a matriz de \emph{reed
        switches} emitiria para o mesmo lance. Nenhuma linha da camada Python
        precisou distinguir os dois casos.
  \item \textbf{As linhas \texttt{@entry} não contaminam o canal de eventos.}
        A verificação final confirma que um leitor que só entenda
        \texttt{casa:estado} continua funcionando --- a extensão do protocolo é
        retrocompatível.
\end{itemize}

\subsection{Rastreabilidade requisito \texorpdfstring{$\times$}{x} teste}
\label{sec:rastreabilidade}

A Tabela~\ref{tab:rastreabilidade} atualiza o mapeamento apresentado na semana
1 (Seção 8.2 do primeiro relatório) com o que esta semana acrescentou, e
registra o que ainda falta validar.

\begin{table}[htbp]
\centering
\caption{Rastreabilidade entre requisitos e evidências, ao fim da semana 3.}
\label{tab:rastreabilidade}
\begin{tabular}{@{}L{0.08\textwidth}L{0.48\textwidth}L{0.34\textwidth}@{}}
\toprule
\textbf{Req.} & \textbf{Evidência} & \textbf{Estado} \\
\midrule
RF1 & \emph{Reed switches}: \texttt{test\_stockfish\_loop} (eventos de sensor
      $\to$ lances) e mock visual. Hardware: 2 fileiras montadas &
      Software \checkmark{} · hardware parcial \\
    & Teclado: 30 verificações de \texttt{test\_keypad\_layer} &
      Não se aplica --- o lance é informado, não detectado \\
RF2 & \texttt{test\_stockfish\_loop} (\emph{engine} falsa, \emph{bestmove}
      aplicado) e \texttt{test\_lichess} / \texttt{test\_challenge} contra o
      servidor falso da \emph{Board API} & \checkmark{} \\
RF3 & Mock com GUI (\texttt{make mock}); destaques de destino e barra de
      status verificados também pela camada de teclado (linhas
      \texttt{@entry}) & Visual \checkmark{} \\
RNF1 & Precisão: 100\% dos eventos digitados produziram o evento esperado
       (30/30). Latência $<$ 200\,ms & Precisão \checkmark{} ·
       \textbf{latência pendente} \\
RNF2 & Instrumentação de \emph{timestamps} na GUI & \textbf{Pendente} \\
RNF3 & Ressincronização (\texttt{0-9}, \texttt{0-0}) e recusa local de eventos
       impossíveis, verificadas na suíte & Software \checkmark{} ·
       \emph{debouncing} em hardware pendente \\
\bottomrule
\end{tabular}
\end{table}

As duas pendências são de medição, não de implementação, e concentram o
trabalho experimental da semana 4: RNF1 e RNF2 exigem instrumentar o caminho
completo (tecla ou sensor $\to$ lance aplicado $\to$ quadro renderizado) no
Raspberry Pi, e RNF3 depende das fileiras montadas para exercitar o
\emph{debouncing} com esbarrões reais.

% ===========================================================================
\section{Documentação e Diagramas}
% ===========================================================================

A documentação foi revisada e ampliada em cerca de 380 linhas no
\texttt{README}, com seções novas para o menu de configuração, as opções
durante a partida, o processo em C e suas duas camadas, o protocolo
\texttt{@entry} e o roteiro de conferência da fiação do teclado (incluindo uma
tabela de sintoma $\to$ causa provável). Todo módulo novo, em C e em Python,
traz cabeçalho explicando a decisão de projeto que ele materializa, e não
apenas o que ele faz.

Quanto aos diagramas, o fluxograma da semana 2 (\texttt{arq\_cpp}) descrevia o
laço de varredura herdado da prova de conceito em Arduino e não corresponde
mais ao código. Ele permanece no repositório, referenciado pelo relatório
daquela semana, mas foi substituído aqui por dois diagramas novos --- mantidos
como fonte em Mermaid (\texttt{docs/diagramas/*.mmd}) e renderizados por
\texttt{mmdc}:

\begin{itemize}
  \item \texttt{arq\_c.mmd} --- arquitetura modular do processo em C, mostrando
        as duas camadas convergindo para o mesmo módulo de IPC
        (Figura~\ref{fig:arq_c.png});
  \item \texttt{fluxo\_keypad.mmd} --- máquina de digitação de um lance,
        incluindo comandos e recusa local (Figura~\ref{fig:fluxo_keypad.png}).
\end{itemize}

Este relatório passou a ser escrito diretamente em \LaTeX{}
(\texttt{docs/Relatorio-3.tex}), em vez de Markdown convertido por Pandoc. O
alvo \texttt{make pdf3} gera os diagramas e compila o documento com
\texttt{xelatex}, tudo dentro do ambiente reprodutível do Nix.

% ===========================================================================
\section{Estruturação do Relatório Final}
% ===========================================================================

O relatório final consolidará os relatórios semanais na estrutura abaixo,
herdada do relatório da semana 1 e ajustada ao que o projeto se tornou. A
coluna de estado indica o que já existe escrito e o que falta.

\begin{table}[htbp]
\centering
\caption{Estrutura prevista para o relatório final.}
\begin{tabular}{@{}L{0.07\textwidth}L{0.46\textwidth}L{0.37\textwidth}@{}}
\toprule
\textbf{Seção} & \textbf{Conteúdo} & \textbf{Estado} \\
\midrule
1 & Motivação e visão geral & Semana 1, revisar \\
2 & Objetivos & Semana 1, revisar o escopo de hardware \\
3 & Requisitos funcionais e não funcionais & Semana 1 (RF1--RF3, RNF1--RNF3); revisar à luz do plano B \\
4 & Ferramentas utilizadas & Atualizar: C no lugar de C++/Arduino, teclado 4$\times$4 \\
5 & Metodologia de desenvolvimento & Semana 1; acrescentar o plano B como decisão de risco \\
6 & Arquitetura física & Semana 1; atualizar com o teclado e a placa efetivamente usada \\
7 & Arquitetura de software (estática e comportamental) & Semana 1 + diagramas novos da semana 3 \\
8 & Implementação: camada Python & Semanas 1--3 \\
9 & Implementação: processo em C e camadas de entrada & Semana 3 (Seções 1.1--1.3) \\
10 & Interface com o usuário & Semana 3 (Seções 1.4--1.5) \\
11 & Hardware: tabuleiro e teclado & Semanas 2--3, pendente o fechamento da montagem \\
12 & Testes: casos planejados, resultados e rastreabilidade & Semana 3 (Seção 3); falta a medição de latência \\
13 & Revisão por pares e evolução do projeto & Semana 2 (Issues nos repositórios dos Grupos P e S) \\
14 & Conclusões, dificuldades e lições & A escrever na semana 4 \\
15 & Referências (ABNT) & Semana 1; acrescentar as citações no texto \\
\bottomrule
\end{tabular}
\end{table}

O limite de 20 páginas de conteúdo (sem referências e anexos) é a restrição
que orienta a consolidação: os relatórios semanais somados excedem esse total,
de modo que as seções de implementação serão condensadas em torno das decisões
de projeto, e os detalhes de operação --- opções de linha de comando, tabelas
de teclas, protocolo IPC --- ficam no \texttt{README}, referenciado a partir do
relatório.

% ===========================================================================
\section{Entrega e Próximos Passos}
% ===========================================================================

A terceira \emph{release} no GitHub marca o estado descrito aqui: processo em C
modular com duas camadas de entrada, interface de configuração e de ações,
suíte com quatro conjuntos de testes e documentação revisada.

Para a semana 4:

\begin{itemize}
  \item medir a latência ponta a ponta (tecla ou sensor até o lance aplicado na
        \emph{engine}), o requisito de $<$ 200\,ms definido na semana 1 --- é a
        evidência experimental que ainda falta;
  \item fechar o que for possível da montagem do tabuleiro e medir a taxa de
        acerto da detecção nas fileiras já montadas;
  \item teste de aceitação: partidas completas contra o Stockfish e contra a IA
        do Lichess pelo teclado, do início ao xeque-mate;
  \item consolidar o relatório final na estrutura da seção anterior, dentro do
        limite de 20 páginas, com as citações no texto em norma ABNT;
  \item registrar em \texttt{docs/figuras/} o material fotográfico da bancada e
        da montagem, como evidência das validações de hardware descritas na
        Seção~2;
  \item recolher as \emph{branches} \texttt{plano\_b} e \texttt{interface}, já
        integradas por \emph{pull request}, deixando apenas \texttt{main} na
        entrega final.
\end{itemize}

\end{document}
